% Tapis de Quillen, with Quillen's offprint and preprint (cote 162-5, typed part): archive cote 162-5, site page 2 (PDF page 3).
% typescript. Draft compared: none.
% AI-transcribed by Claude Opus 5.5 (claude-opus-5-5, adaptive thinking, effort high), Message Batches API, batch t162_full (msgbatch_01UAh3zyWPd454AeDkhnmQ33), 2026-11-13; not yet checked by eye.
\scanpage{162-5/2}
\mspage{\hand{1}} \margin{notes du \uline{10.9.68}}

Tapis de Quillen.

I \quad \uline{Relations entre catégories et ensembles semi-simplic\unsure{ia}ux.}

A toute catégorie C, \ on associe un ensemble semi-simplicial S(C),
trouvant ainsi un foncteur pleinement fidèle \xout{de Cat}
\[ S;\ (\mathrm{Cat}) \longrightarrow (\mathrm{Ssimpl}) \quad . \]
\xout{Si C est associé à} Les systêmes locaux d'ensemble sur SC correspondent
aux foncteurs sur C qui transforment toute flêche en isomorphisme (i.e.
qu\unsure{i} se factorisent par le groupoïde associé à \ C). Les H$^{i}$ sur SC d'un tel
système local (H$^{o}$ pour ens, H$^{1}$ pour groupes, H$^{i}$ quelconques pour groupes
abéliens) s'interprètent en termes des foncteurs $\uline{\mathrm{lim}}^{(i)}$ dérigés de $\uline{\mathrm{lim}}$ ,
ou si on préfère, des H$^{i}$( du \uline{topos} C\hand{\unsure{$\widetilde{\ }$}}. On voit ainsi à quelle condition
un foncteur \ C $\longrightarrow$ C' induit un\xout{e} \ homotopisme \ SC $\longrightarrow$ SC' : \ en vertu du
critère cohomologique de Artin-Mazur, il f et s que \ pour tout systême
de coéfficients \ F'sur C', l'homomorphisme naturel $\underset{C'}{\underleftarrow{\mathrm{lim}}}^{(i)}F' \longrightarrow \underset{C}{\underleftarrow{\mathrm{lim}}}^{(i\ )}F$
soit un isomorphisme (pour les i pour lesquels cela a un sens). \xout{\unsure{Exemples}}

A C on peut asso\unsure{ci}er le topos $\hand{\widetilde{C}}$ , qui varie de façon \uline{covariante} avec C
% NOTE: right edge of the scan cuts off the ends of this and the following lines; the visible letters are transcribed
(NB le foncteur C$\longmapsto$$\hand{\widetilde{C}}$ \ n'a plus rien de pleinement fidèle, semble-t\unsure{-}il ??)
Les systêmes de coéfficients \inserted{ensemblistes} sur C ( \inserted{déf} = \ les foncteurs \inserted{$C^{o} \longrightarrow$ Ens} transformant isomo\unsure{r-}%
phismes en isomorphismes) correspondent aux faisceaux localement constan\unsure{ts}
i.e. les objets localement constants de $\hand{\widetilde{C}}$, définis intrinsèquement en te\unsure{rmes}
de $\hand{\widetilde{C}}$ . \ Ainsi, le fait pour un foncteur F:C$\longrightarrow$C' d'induire un\xout{e équivalence}
homotopisme S(C)$\longrightarrow$S(C') \ ne \ dépend que du morphisme de topos $\hand{\widetilde{F}}$: $\hand{\widetilde{C}} \longrightarrow \hand{\widetilde{C'}}$
induit, et signifie que pour tout faisceau localement constant F' sur C'
c. sur $\hand{\widetilde{C'}}$ , \ les applic\unsure{a}tions induites \ H$^{i}$($\hand{\widetilde{C'}}$,F') $\longrightarrow$ H$^{i}$($\hand{\widetilde{C}}$,$\hand{\widetilde{F}}^{*}$(F'))
sont des isomorphismes (pour les i pour lesquels cela a un\xout{\unsure{es}} sens).

\margin{2}
